% ============================================================
% 5_Analysis_and_Discussion.tex
% Real data: Finding 1 (FD values, WA), Finding 2 (prosody gap, 3-seed/2-seed),
%   Finding 3 (v5 augmentation), Finding 4 (C-BESD layer weights),
%   APC (540-sample mean±std). [待B1] markers remain for FAU/IEMOCAP layer weights.
% ============================================================

\section{Analysis and Discussion}
\label{sec:analysis_discussion}

The 63-experiment matrix reveals four interlocking regularities that,
collectively, argue for distribution awareness as a necessary design principle
in children's SER.

\subsection{Finding 1: FD--Accuracy Monotonic Relationship}

Across enhancement (C1--C4), age (child vs.\ adult), and style (acted vs.\
spontaneous) dimensions, Fr\'{e}chet Distance between corpus-level WavLM
feature distributions exhibits a strong monotonic negative correlation with
classification accuracy. Table~\ref{tab:fd_wa} summarizes this relationship
using the best measured accuracy for each condition.

\begin{table}[h]
\centering
\caption{FD--Accuracy correspondence across distribution-shift dimensions. All FD values from canonical AutoDL GPU computation; accuracy from E1/E3/E4/X1 best pooling.}
\label{tab:fd_wa}
\begin{tabular}{lccc}
\toprule
\textbf{Shift Dimension} & \textbf{FD} & \textbf{Best WA (\%)} & \textbf{Experiment} \\
\midrule
Same corpus (in-domain)       & 0.00  & 94.97 & E1-02 (C-BESD SA, 3-seed) \\
Age shift (adult in-domain)   & 7.20  & 75.96 & E1-08 (IEMOCAP SA, 2-seed) \\
Style shift (spont.\ in-dom.) & 8.50  & 66.46 & E1-05 (FAU SA, 3-seed) \\
Style shift (zero-shot)       & 8.50  & 19.56 & E3 (C$\rightarrow$F) \\
Enhancement shift (extreme)   & 11.99 & 46.50 & E4-04 (C-BESD extreme, v5) \\
Language shift (EN$\rightarrow$TE) & 16.48 & 19.68 & X1 (cross-lingual) \\
\bottomrule
\end{tabular}
\end{table}

\noindent The monotonicity of this relationship---every increase in FD
predicts a decrease in WA---suggests that FD is a reliable diagnostic for
anticipating cross-corpus performance degradation \textit{before} training. The
zero-shot style-shift condition (FD=8.50, WA=19.56\%) and the in-domain
style-shift condition (FD=8.50, WA=66.46\%) share the same FD but differ by
46.9pp in accuracy, underscoring that FD measures corpus-level distribution
distance while in-domain training can partially overcome this gap through
corpus-specific adaptation. This has practical implications for corpus design
and model selection in child SER: if FD between the training corpus and target
deployment population exceeds a threshold, in-domain training on a matched
corpus should be prioritized over zero-shot transfer.

\subsection{Finding 2: Population-Dependent Prosody Pooling Efficacy}

Table~\ref{tab:prosody_gap} compares Self-Attention (SA) and Prosody-Guided
(PG) pooling across populations. On children's speech, SA and PG achieve
near-identical performance: $\Delta$(SA$-$PG) $=-0.13$pp on C-BESD (PG
marginally better, not significant at $p>0.05$) and $+1.31$pp on FAU
(SA marginally better). On adult speech (IEMOCAP), SA outperforms PG by
$+9.73$pp---a large and meaningful gap.

\begin{table}[h]
\centering
\caption{Prosody-Guided vs.\ Self-Attention gap by population. C-BESD/FAU: 3-seed; IEMOCAP: 2-seed.}
\label{tab:prosody_gap}
\begin{tabular}{lccc}
\toprule
\textbf{Dataset} & \textbf{Population} & \textbf{$\Delta$WA (SA $-$ PG)} & \textbf{Direction} \\
\midrule
C-BESD & Children (acted) & $-$0.13pp & PG $\approx$ SA (n.s.) \\
FAU Aibo & Children (spontaneous) & +1.31pp & SA $\approx$ PG (marginal) \\
IEMOCAP & Adults (acted) & +9.73pp & SA $\gg$ PG (large gap) \\
\bottomrule
\end{tabular}
\end{table}

\noindent This pattern reverses the common intuition that prosody priors are
more informative for children. Instead, the data show that frame-level F0 and
energy priors are \emph{neutral} for children's SER---neither helping nor
hurting---but are \emph{actively harmful} for adult SER. We interpret this
through the lens of SSL representational sufficiency: WavLM representations
already encode prosodic information at mid-to-upper layers (Finding~4).
Adding explicit prosody features to the attention computation introduces
redundant information for adult speech (where prosodic patterns are more
stereotyped and already captured by SSL) but is harmless for children's speech
(where higher F0 variability and distinct energy contours make the redundant
signal less interfering). The practical implication is clear: \textbf{for
children's SER, prosody guidance is safe but unnecessary; for adult SER, it is
harmful and should be avoided.}

\subsection{Finding 3: Non-Transferability of Adult Augmentation}

Experiments E4-01--E4-04 (C-BESD, v5 protocol) provide clear evidence that
adult-derived augmentation strategies do not transfer to children. Adding adult
IEMOCAP speech to C-BESD training data degrades WA by 30.63pp
(81.24\% $\rightarrow$ 50.61\%), making it the second-most harmful condition
after extreme augmentation ($-$34.74pp). Child-matched AWGN is less harmful
($-$21.35pp), confirming that the distributional mismatch from adult speech is
more damaging than controlled noise injection. The FD--WA relationship holds
across all four conditions: FD 0.00 $\rightarrow$ 81.24\%, FD 8.71
$\rightarrow$ 59.89\%, FD 9.87 $\rightarrow$ 50.61\%, FD 11.99 $\rightarrow$
46.50\%.

This finding has a clear practical implication: \textbf{do not augment
children's SER training data with adult speech.} The common practice of pooling
all available emotional speech data for pre-training is actively harmful when
the target population is children. Furthermore, FD serves as a reliable
pre-augmentation diagnostic: the FD between the augmentation source and target
corpus predicts the expected degradation.

\subsection{Finding 4: Universality of Mid-to-Upper Layer Preference}

Across all three datasets, the learned layer fusion weights consistently assign
highest values to mid-to-upper layers (L7--L9, 1-based), with low but positive
weights on lower layers and near-zero weights on the final layer.
Table~\ref{tab:layer_weights} reports the learned weight distribution from the
Exp2 (C-BESD, Prosody-Guided) checkpoint. FAU and IEMOCAP layer weights remain
pending B1 experiments.

\begin{table}[h]
\centering
\caption{Learned layer fusion weights from Exp2 checkpoint (C-BESD, Prosody-Guided). Theoretical max entropy = $\log 12 \approx 2.485$.}
\label{tab:layer_weights}
\begin{tabular}{lccc}
\toprule
\textbf{Dataset} & \textbf{Argmax Layer (1-based)} & \textbf{Entropy} & \textbf{Weight Range} \\
\midrule
C-BESD & L9 & 2.484 & 0.04--0.16 \\
FAU Aibo & --$^\dagger$ & --$^\dagger$ & --$^\dagger$ \\
IEMOCAP & --$^\dagger$ & --$^\dagger$ & --$^\dagger$ \\
\bottomrule
\end{tabular}
\end{table}

\noindent The near-maximum entropy (2.484 out of 2.485) indicates that all 12
layers contribute near-uniformly, yet the systematic peak at L9 (argmax weight
$\approx$0.091) confirms that emotion-discriminative information concentrates in
the upper-middle layers. This pattern---near-uniform distribution with a
consistent upper-middle peak---is consistent with prior findings that WavLM's
intermediate layers (7--10) encode the most transferable speech representations.
The argmax position (L9, 1-based) aligns with layer-wise probing results
showing that paralinguistic information peaks in WavLM layers 8--10.

\noindent$^\dagger$Pending B1 layer fusion experiments on FAU and IEMOCAP.

\subsection{Attention-Prosody Correlation (APC)}

For models using Prosody-Guided pooling, we compute the Attention-Prosody
Correlation (APC$_{\text{wav}}$)---the Pearson correlation between frame-level
attention weights and frame-level acoustic energy---over all 540 C-BESD test
samples. APC$_{\text{wav}} = 0.7406 \pm 0.1087$ (mean$\pm$std), with
APC$_{\Delta} = -0.0450 \pm 0.0891$ (the difference between Prosody-Guided and
Self-Attention APC). The negative APC$_{\Delta}$ indicates that Prosody-Guided
pooling \textit{reduces} redundant attention--acoustic correlation relative to
pure self-attention: by injecting explicit prosody information into the
attention computation, the model learns to rely \textit{less} on the implicit
prosody correlations already present in SSL features. This mechanism explains
why Prosody-Guided pooling does not harm children's SER (Finding~2): the
prosody prior serves as a regularization signal that decorrelates attention
from raw acoustic energy, counterbalancing any redundancy it introduces.

\subsection{Finding 5: Feature Space Visualization via t-SNE}

To qualitatively assess how the distribution-driven pipeline transforms the
feature space, we apply t-SNE~\cite{van2008visualizing} to the penultimate
classifier layer (128-dim) before and after the full pipeline. Following
Neumann \& Vu~\cite{neumann2019cross}, who visualized the last hidden layer
of an attentive CNN for SER, we extract features from the layer immediately
preceding the final linear projection. For the ``before'' condition, we use
raw frozen WavLM last-layer output (768-dim, mean-pooled over time) without
any pipeline modules. For the ``after'' condition, we use the full
distribution-driven pipeline (WavLM $\rightarrow$ LayerFusion $\rightarrow$
Adapter $\rightarrow$ Pooling $\rightarrow$ SEMLP penultimate 128-dim).
All t-SNE projections use perplexity 30 with stratified per-class sampling
(2,000 samples per dataset).

Figure~\ref{fig:tsne} shows the 3$\times$2 t-SNE grid across C-BESD
(children acted, 6-class), FAU Aibo (children spontaneous, 4-class), and
IEMOCAP (adults acted, 4-class). Three patterns emerge:

\begin{enumerate}
    \item \textbf{Clear separation gain on acted children's speech.}
    For C-BESD (WA=91.87\%), the before panel shows partially overlapping
    clusters, while the after panel reveals six well-separated emotion
    clusters, consistent with the strong quantitative performance.

    \item \textbf{Spontaneous speech resists clean separation.}
    For FAU Aibo (WA=67.05\%, UAR=42.92\%), the ``neutral'' class remains
    diffusely distributed across all regions even after pipeline processing.
    The large WA--UAR gap ($\sim$24pp) is visually corroborated: the model
    achieves moderate WA by capturing majority-class boundaries but fails
    to form compact per-class clusters, particularly for under-represented
    emotions. This visual evidence reinforces that spontaneous children's
    speech is the hardest condition and that distribution awareness alone
    cannot fully compensate for within-corpus class imbalance.

    \item \textbf{Adult speech benefits from pipeline structure.}
    For IEMOCAP (WA=63.76\%), the pipeline transforms dispersed raw WavLM
    features into moderately structured clusters, though separation remains
    weaker than C-BESD---consistent with the cross-population FD gap.
\end{enumerate}

\noindent To our knowledge, this is the first t-SNE visualization of
children's speech emotion features under a distribution-driven pipeline,
addressing a noted gap in the SER visualization literature~\cite{kanwal2022speech}.

\subsection{Limitations}

We identify four limitations. First, our experimental coverage is bounded by
available GPU resources (67 experiments on RTX 4090D); broader architectures
(e.g., HuBERT, data2vec) and larger child corpora (MyST, KidsTALC) remain for
future work. Second, C-BESD child age metadata is unavailable, preventing
fine-grained developmental analysis. Third, our FD measurements use WavLM
features---the metric is representation-dependent, and FD values may differ
under different SSL backbones. Fourth, all three datasets are from WEIRD
(Western, Educated, Industrialized, Rich, Democratic) populations;
cross-cultural child SER remains unexplored.
